% Section 5.4, from lectures/L05/README.md: what you should be able to do, questions to test
% yourself, and the next lecture.

\section{Review}
\label{c5:sec:review}

\needspace{6\baselineskip}
\subsection*{What you should be able to do now}
\begin{itemize}
  \item Say what a prescaler does, name the five ratios, and say which steps between them are not a
    factor of eight.
  \item Configure Timer1 for CTC mode at a chosen prescaler, from the datasheet's bit names.
  \item Turn a wanted frequency into a prescaler and a compare value, by hand.
  \item Say why the compare value is one less than the tick count, and what going wrong there costs.
  \item State the fastest and slowest frequencies a given timer can reach, and say what to do about
    a frequency below the slowest.
  \item Explain why an 8-bit timer is so much worse at hitting a frequency exactly than a 16-bit
    one, in terms of ticks rather than of bits.
  \item Extend a timer in software so that one interrupt drives several events at different rates.
  \item Predict the number of cycles between two compare-match interrupts, and say when that number
    is the same every time and when it is not.
\end{itemize}

\needspace{6\baselineskip}
\subsection*{Questions to test yourself}
\begin{enumerate}
  \item Why is a delay loop a worse way to wait than a timer, given that both work?
  \item The prescaler ratios are 1, 8, 64, 256 and 1024. Which steps are the odd ones, and what
    does that mean when a frequency does not fit at one prescaler?
  \item You want 1\,kHz from Timer1 at 16\,MHz. Give the prescaler and the compare value, and say
    whether the result is exact.
  \item You want the same 1\,kHz from Timer0. Give the prescaler and compare value, and say why
    they are different numbers for the same frequency.
  \item An 8-bit timer at 16\,MHz cannot produce 1\,Hz. Say why, in one sentence, and say what you
    would do instead.
  \item Two compare-match interrupts are 5333 cycles apart on average, but no individual gap is
    5333. What is going on, and is the timer at fault?
\end{enumerate}

\needspace{6\baselineskip}
\subsection*{Looking ahead}
Chapter~6 finishes the library, and connects it to the language everything else is written in:
\begin{itemize}
  \item The watchdog: a timer whose job is to notice that your program has stopped.
  \item Timed write sequences, and hardware that refuses to be configured carelessly.
  \item Sleep modes, and what is still running in each.
  \item Calling your assembly from C, and C from your assembly, which is where the calling contract
    you have been following since Chapter~1 finally has another party to it.
\end{itemize}
